\documentclass[11pt]{article}
\usepackage[margin=1in]{geometry}
\usepackage{amsmath,amssymb,amsthm}
\usepackage{enumitem}
\usepackage{booktabs}
\usepackage{microtype}
\usepackage{hyperref}
\hypersetup{colorlinks=true,urlcolor=blue,linkcolor=black,citecolor=black}

\setlength{\parskip}{4pt plus 1pt}
\setlength{\parindent}{0pt}

\newtheorem{thm}{Theorem}
\newtheorem{lem}[thm]{Lemma}
\newtheorem{prop}[thm]{Proposition}
\newtheorem{cor}[thm]{Corollary}
\newtheorem{conj}[thm]{Conjecture}
\theoremstyle{definition}
\newtheorem{defn}[thm]{Definition}

\title{Step-1 kill statistics and asymptotic locality:\\
       \large empirical evidence for survivor fraction $\to 0$}
\author{Jonathan Valenzuela \\ \small{(satellite for the locality program)}}
\date{}

\begin{document}

\maketitle

\begin{abstract}
At $n$ external legs, the layer-$0$ (``Step-1'') hidden-zero kill mechanism
forces $a_M = 0$ for the coefficient of every chord multiset $M$ of size
$n-3$ that satisfies the killability conditions (K1)+(K2) at some cyclic
zone. Through $n=9$ we report the exact count of surviving (non-killable)
non-triangulation multisets at each $n$, the resulting kill rate, and the
cyclic-orbit structure of the survivors. The data show: (i) the kill rate
on non-triangulations exceeds $99.6\%$ for every $n \ge 7$ measured;
(ii) the kill rate is monotonically increasing in $n$; (iii) at $n=9$
already $99.888\%$ of all non-triangulation multisets are killed by
Step-1 alone. We state the precise asymptotic conjecture (Item 4 of the
master plan) and identify what a proof would need to establish.
This satellite is companion to the local-survival lemma of
\verb|paper/step-one-doesnt-kill-triangulations/|, which closes the
matching combinatorial side (Item 2).
\end{abstract}

%--------------------------------------------------------------------------
\section{Setup}
%--------------------------------------------------------------------------

We adopt the conventions of the root README: $n$ vertices labelled
$1, \dots, n$ cyclically; planar Mandelstam invariants
$X_{ij} = (p_i + \cdots + p_{j-1})^2$ for chord pairs $(i, j)$; rational
ansatz
\[
   B \;=\; \sum_M \frac{a_M}{\prod_{c \in M} X_c}
\]
where $M$ runs over size-$(n-3)$ chord multisets. The cyclic
hidden-zero zone $\mathcal Z_{r, r+2}$ has special chord
$X_{r,r+2}$, bare chord $X_{r+1,r-1} = -X_{r,r+2}$, and substitutions
$X_{r+1, k} = X_{r,k} - X_{r,r+2}$ for $k$ in cyclic range.

$M \in \mathcal K_{r,r+2}$ (``killable at zone $r$'') iff
\begin{itemize}[leftmargin=*,topsep=0pt,itemsep=2pt]
\item[\rm (K1)] $M$ contains neither the special nor the bare;
\item[\rm (K2)] for every (companion, substitute) pair $(X_{r,k}, X_{r+1,k})$,
$M$ does not contain both members.
\end{itemize}
A multiset $M$ is a \emph{Step-1 survivor} iff it is in no
$\mathcal K_{r, r+2}$ for any $r$. By the local-survival lemma of
\cite{tri-survives}, every triangulation $T$ is a Step-1 survivor; the
non-triangulation Step-1 survivors are the focus of this note.

%--------------------------------------------------------------------------
\section{The data}
%--------------------------------------------------------------------------

The script \verb|computations/step1_layer0_kill/kill_enumeration/step1_uncaught.py|
enumerates all size-$(n-3)$ multisets, applies the (K1)+(K2) test at
every cyclic zone, and reports the survivors. The output through
$n = 9$:

\begin{center}
\begin{tabular}{rrrrrr}
\toprule
$n$ & total multisets & triangulations $C_{n-2}$
       & non-tri total & non-tri survivors & non-tri kill \% \\
\midrule
4 &       2 &     2 &       0 &     0 & --- \\
5 &      15 &     5 &      10 &     0 & 100\,\% \\
6 &     165 &    14 &     151 &     0 & 100\,\% \\
7 &   2 380 &    42 &   2 338 &     7 & 99.700\,\% \\
8 &  42 504 &   132 &  42 372 &   100 & 99.764\,\% \\
9 & 906 192 &   429 & 905 763 & 1 011 & 99.888\,\% \\
\bottomrule
\end{tabular}
\end{center}

Triangulation counts match the Catalan numbers $C_{n-2}$ exactly,
consistent with the local-survival lemma of \cite{tri-survives}. The
total multiset count at $n$ is
$\binom{d+n-4}{n-3}$ where $d = n(n-3)/2$.

\paragraph{Kill rate is monotonically increasing.} For every $n \ge 7$
in the data, the kill rate strictly exceeds the rate at $n - 1$:
$99.700\% \to 99.764\% \to 99.888\%$. The headline number at $n = 9$
is that fewer than $0.12\%$ of all non-triangulation multisets escape
Step-1.

%--------------------------------------------------------------------------
\section{Cyclic orbit structure of the survivors}
%--------------------------------------------------------------------------

The survivors at each $n$ are closed under the cyclic-rotation action
$v \mapsto v + 1 \pmod{n}$ on chord endpoints, so they break into
$\mathbb Z_n$-orbits of size $\le n$.

\paragraph{$\boldsymbol{n=7}$.} All 7 survivors form a single
$\mathbb Z_7$-orbit (the ``fish'' family); they are the cyclic shifts of
$M_1 = \{(1,3), (1,6), (2,4), (4,6)\}$. See
\verb|paper/Laurent series for hard kills/cascade_n7.tex| for the kill.

\paragraph{$\boldsymbol{n=8}$.} The 100 survivors break into multiple
$\mathbb Z_8$-orbits (orbit-shape detail in
\verb|computations/step1_layer0_kill/kill_enumeration/outputs/|). The
companion script
\verb|computations/step3_laurent/cascade_n8/cascade_kill_n8.py| verifies
\textbf{every one} of these 100 survivors dies at Step~3 via a depth-$1$
Laurent cascade --- 90 of them at order $X_*^{-3}$ (leading + 1) and 10
at order $X_*^{-4}$ (leading + 2), with all cousins step-$1$ killable at
one neighbouring zone. So at $n=8$, the cumulative kill rate of
\textbf{Step-1 + Step-3} on non-triangulations is $100\%$.

\paragraph{$\boldsymbol{n=9}$.} 1011 survivors. Cyclic-orbit
decomposition is in the script's PDF output. The Step-3 cascade
verification at $n=9$ has not yet been run.

%--------------------------------------------------------------------------
\section{Asymptotic conjecture (Item 4)}
%--------------------------------------------------------------------------

Let $S(n)$ denote the number of non-triangulation Step-1 survivors at
$n$, and $T(n)$ the total number of non-triangulation size-$(n-3)$
multisets ($T(n) = \binom{d + n - 4}{n-3} - C_{n-2}$). The data
$(S(n), T(n))$ for $n = 5, \dots, 9$:

\begin{center}
\begin{tabular}{rrrcc}
\toprule
$n$ & $S(n)$ & $T(n)$ & $S(n)/T(n)$ & $S(n+1)/S(n)$ \\
\midrule
5 &     0 &      10 &  0\,\%       &   ---     \\
6 &     0 &     151 &  0\,\%       &   ---     \\
7 &     7 &   2 338 &  0.300\,\%   &   ---     \\
8 &   100 &  42 372 &  0.236\,\%   & $\approx 14.3$ \\
9 & 1 011 & 905 763 &  0.112\,\%   & $\approx 10.1$ \\
\bottomrule
\end{tabular}
\end{center}

The ratio $S(n)/T(n)$ is roughly halving per increment of $n$, while
$S(n+1)/S(n)$ is decreasing (suggesting sub-exponential survivor
growth). The total grows superexponentially.

\begin{conj}[Asymptotic locality from Step-1 alone]\label{conj:asympt}
There exist constants $C, \alpha > 0$ such that for every $n \ge 7$,
\[
   \frac{S(n)}{T(n)} \;\le\; \frac{C}{n^{\alpha}}
   \quad\text{(polynomial decay)}\qquad\text{or}\qquad
   \frac{S(n)}{T(n)} \;\le\; C \cdot 2^{-n}
   \quad\text{(exponential decay).}
\]
The data favour the exponential form with rate $\sim 2^{-n}$ but the
polynomial form is not ruled out by $n \le 9$.
\end{conj}

A proof would need to:
\begin{enumerate}[leftmargin=*,topsep=0pt,itemsep=2pt]
\item identify the structural family of multisets that satisfies neither
       (K1) nor (K2) at every zone (the survivors);
\item bound the cardinality of that family by a closed-form expression
       in $n$;
\item compare to $T(n)$, whose growth is known.
\end{enumerate}
The cyclic-orbit structure reduces (1)--(2) to characterising one
representative per orbit, then multiplying by orbit size $\le n$.
At $n = 7$ all 7 survivors are one orbit; at $n = 8$ the orbit count is
the natural quantity to track in the proof. We do not pursue this here.

\paragraph{What this conjecture is \emph{not} needed for.} The all-$n$
locality theorem does \emph{not} depend on
Conjecture~\ref{conj:asympt}. What it depends on is the per-survivor
Laurent cascade (Item 10 of the master plan): every Step-1 survivor
dies via a depth-1 cascade with prerequisites all of layer-0 type.
Verified at $n = 7$ symbolically and at $n = 8$ in
\verb|computations/step3_laurent/cascade_n8/|. The asymptotic
conjecture here is supportive evidence, not a load-bearing component.

%--------------------------------------------------------------------------
\section{What the satellite Step-1 ecosystem now looks like}
%--------------------------------------------------------------------------

The two satellites that close out the Step-1 analysis:

\begin{itemize}[leftmargin=*,topsep=2pt,itemsep=4pt]
\item \verb|paper/step-one-doesnt-kill-triangulations/| --- proves the
       \emph{local} side: no triangulation is ever Step-1 killable.
       Closes Item 2 of the master plan.
\item \verb|paper/step-one-kill-statistics/| (this note) --- documents
       the \emph{global} kill rate: which fraction of non-triangulations
       does Step-1 catch, and how does that fraction grow with $n$?
       Backs Item 4 of the master plan.
\end{itemize}

Together with the homogeneity-decomposition lemma
(\verb|paper/B=0->B_i = 0/|, Item 1) and the Laurent-cascade analysis
(\verb|paper/Laurent series for hard kills/|, Item 10), Step-1 is
fully understood up to the asymptotic question above.

%--------------------------------------------------------------------------
\begin{thebibliography}{9}
\bibitem{tri-survives}
J.\ Valenzuela, \emph{Layer-0 hidden-zero kills no triangulation},
\verb|paper/step-one-doesnt-kill-triangulations/triangulation_layer0.pdf|.

\bibitem{cascade-n7}
J.\ Valenzuela, \emph{The Laurent-tail kill at $n=7$},
\verb|paper/Laurent series for hard kills/cascade_n7.pdf|.

\bibitem{rodina}
A.\ Rodina, \emph{Hidden zeros are equivalent to enhanced ultraviolet
scaling and lead to unique amplitudes in $\mathrm{Tr}(\phi^3)$ theory},
arXiv:2406.04234.
\end{thebibliography}

\end{document}
